\chapter{روش‌شناسی}
\label{ch:methodology}

\section{طراحی مطالعه}
این پژوهش یک مطالعۀ گذشته‌نگر مبتنی بر داده‌های ساخت‌یافتۀ \lr{MIMIC\text{-}IV v3.1} است. هدف، پیش‌بینی نشانگر جانشین امکان‌پذیری ترخیص از \lr{ICU} با استفاده از داده‌های \pnum{24} ساعت نخست بستری است. تحلیل روی کوهورت کامل واجد شرایط انجام شد و هیچ سقف نمونه‌گیری فعال نبود.

\section{منبع داده و کامل بودن داده خام}
داده‌های خام از مسیر محلی \texttt{\lr{mimic-iv-3.1/}} خوانده شدند. اعتبارسنجی نهایی نشان داد فایل‌های ضروری بیمارستان و \lr{ICU}، از جمله \texttt{\lr{patients.csv.gz}}، \texttt{\lr{admissions.csv.gz}}، \texttt{\lr{diagnoses\_icd.csv.gz}}، \texttt{\lr{labevents.csv.gz}}، \texttt{\lr{icustays.csv.gz}} و \texttt{\lr{chartevents.csv.gz}} موجود بودند. گزارش اعتبارسنجی شامل \pnum{33} ورودی در مانیفست و \pnum{33} فایل موجود بود؛ بنابراین مانیفست کامل تأیید شد.

\section{تعریف کوهورت}
کوهورت شامل بستری نخست بزرگسال در \lr{ICU} پس از معیارهای ورود و خروج مطالعه بود. ردیف‌های خارج از این کوهورت وارد مدل‌سازی نشدند. مجموعۀ نهایی دارای \textbf{\pnum{32,399}} ردیف، \pnum{29} ستون، \pnum{22} ستون ویژگی، \pnum{32,399} بستری یکتای \lr{ICU} و \pnum{32,399} بیمار یکتا بود. همۀ ردیف‌ها پس از معیارهای مطالعه استفاده شدند و مقدار \lr{max\_cohort\_sample} برابر \lr{null} بود.

\section{تعریف برچسب}
برچسب مثبت نشان می‌دهد بیمار طبق تعریف جانشین، پیامد امکان‌پذیر داشته است. سه شرط هم‌زمان لازم بود: زنده ماندن تا ترخیص بیمارستان، مقصد ترخیص مطلوب، و طول اقامت \lr{ICU} کمتر از \pnum{720} ساعت. این تعریف از داده‌های ساخت‌یافته استخراج می‌شود و به‌معنای آمادگی واقعی ترخیص یا تصمیم درمانی نیست.

توزیع نهایی برچسب‌ها چنین بود: \pnum{12,102} نمونه در کلاس \pnum{0} و \pnum{20,297} نمونه در کلاس \pnum{1}. نسبت کلاس مثبت حدود \pnum{62.6}\% بود.

\section{استخراج ویژگی}
ویژگی‌ها از پنجرۀ \pnum{24} ساعت نخست بستری \lr{ICU} استخراج شدند. این ویژگی‌ها شامل خلاصه‌های آزمایشگاهی، علائم حیاتی، و متغیرهای مشتق‌شده بودند. پیش از جایگذاری، میانگین درصد گمشده‌ها \cnum{0.252} و بیشترین درصد گمشده‌ها \cnum{0.998} بود. پس از جایگذاری، تعداد سلول‌های گمشده در دادۀ مدل‌سازی صفر گزارش شد.

\section{توکن‌های علائم}
در کنار ویژگی‌های جدولی، بازنمایی توکنی علائم ساخته شد. خروجی نهایی در فایل \texttt{\lr{data/processed\_full\_cohort/tokens.npy}} ذخیره شد و شکل آن \lr{(32399, 25, 4)} بود. هر نمونه دارای \pnum{25} گام توکنی و \pnum{4} بعد رمزگذاری بود. این توکن‌ها برای مدل \lr{Few-Shot ETHOS} و مدل هیبریدی استفاده شدند.

\begin{figure}[htbp]
\centering
\begin{tikzpicture}[
  node distance=1.2cm,
  box/.style={draw, rounded corners=2pt, align=center, minimum width=3.2cm, minimum height=0.9cm},
  arrow/.style={-Latex, thick}
]
\node[box] (raw) {داده خام \\ \lr{MIMIC-IV v3.1}};
\node[box, left=of raw] (cohort) {معیارهای ورود/خروج \\ بستری نخست بزرگسال};
\node[box, below=of raw] (features) {ویژگی‌های \pnum{24} ساعت نخست \\ آزمایش و علائم حیاتی};
\node[box, right=of raw] (label) {برچسب جانشین \\ ترخیص از \lr{ICU}};
\node[box, below=of features] (tokens) {توکن‌های علائم \\ \lr{(32399,25,4)}};
\node[box, right=of tokens] (models) {مدل‌های جدولی، هیبریدی، \\ چندنمونه‌ای و فدرال};
\node[box, below=of models] (eval) {ارزیابی: \\ \lr{AUROC, AUPRC, Brier, ECE}};
\draw[arrow] (cohort) -- (raw);
\draw[arrow] (raw) -- (features);
\draw[arrow] (raw) -- (label);
\draw[arrow] (features) -- (tokens);
\draw[arrow] (tokens) -- (models);
\draw[arrow] (label) |- (models);
\draw[arrow] (models) -- (eval);
\end{tikzpicture}
\caption{جریان کلی استخراج داده، ساخت توکن‌ها و ارزیابی مدل‌ها}
\label{fig:workflow}
\end{figure}

\section{مدل‌ها}
مدل‌های اصلی عبارت بودند از رگرسیون لجستیک، جنگل تصادفی، \lr{XGBoost}، \lr{LightGBM}، \lr{CatBoost}، خط پایۀ عصبی، \lr{Few-Shot ETHOS}، \lr{Token Hybrid RF}، مدل انباشتی کلاسیک، و \lr{Federated LR Node Ensemble}. مدل انباشتی پیش‌بینی‌های مدل‌های کلاسیک تکمیل‌شده را با یک فرا-یادگیر ترکیب کرد. مدل فدرال با تقسیم گره‌ای شبیه‌سازی‌شده آموزش داده شد و خروجی گره‌ها را تجمیع کرد.

\section{تقسیم داده و معیارهای ارزیابی}
داده به سه بخش آموزش، اعتبارسنجی و آزمون تقسیم شد: \pnum{20,735} ردیف آموزش، \pnum{5,184} ردیف اعتبارسنجی و \pnum{6,480} ردیف آزمون. آستانه‌ها با استفاده از مجموعۀ اعتبارسنجی تنظیم شدند و گزارش نهایی بر مجموعۀ آزمون نگه‌داشته‌شده انجام شد.

معیارهای گزارش‌شده شامل \lr{AUROC}، \lr{AUPRC}، دقت، بازخوانی یا حساسیت، ویژگی، \lr{F1}، امتیاز بریر، \lr{ECE}، عرض از مبدأ و شیب کالیبراسیون، و ماتریس درهم‌ریختگی بودند. استفاده از معیارهای کالیبراسیون برای جلوگیری از تفسیر بیش‌ازحد \lr{AUROC} ضروری است \cite{steyerberg2019clinical,vanCalster2019calibration}.

\section{بازتولیدپذیری و اخلاق}
کد استخراج و مدل‌سازی مسیرهای جداگانۀ \texttt{\lr{data/processed\_full\_cohort/}} و \texttt{\lr{results/full\_cohort/}} را استفاده می‌کند. دادۀ خام \lr{MIMIC\text{-}IV} قابل بازتوزیع نیست و فقط نتایج تجمیعی، کد و پیکربندی قابل اشتراک‌اند. داده‌ها بی‌نام‌سازی شده‌اند، دسترسی نیازمند احراز صلاحیت \lr{PhysioNet} و توافق‌نامه استفاده از داده است، تماس مستقیم با بیمار وجود نداشت، و هیچ تلاشی برای بازشناسایی بیماران انجام نشد. وضعیت معافیت یا تأیید نهاد اخلاق دانشگاهی باید توسط نویسنده با مقررات محلی تأیید شود.
